% Regression fixture: a {tasks} grid is a list in the tagged PDF.
%
% The tasks package sets each task in a box of its own, and a tagged build used
% to record a row as paragraphs nested in a paragraph: veraPDF failed the file
% on ISO 32005 Table 5, "<P> shall not contain <Part>" and "<P> shall not
% contain <P>", once each for the text and for the label of every task after
% the first in its row. Without texlib-tasks.sty this fixture fails 50 checks.
% That package retags the grid as a list. The fixture holds one grid for each
% way that can go wrong, and it is asserted twice, both times in the accessible
% build (examples/manifest.py, run by smoke_test.py --accessible):
%
%   - veraPDF must pass it, which is the nesting;
%   - tagged= counts its lists and list items, which is the structure. A grid
%     that validates but is not a list (every task a loose paragraph) passes the
%     first check and fails the second.
%
% The count is 11 lists and 27 items: nine grids of 23 tasks, the {parts} list
% and the {itemize} of two items each. A grid that came out as one item is 13
% items; a grid that came out as a list inside a one-item list is 20 lists.
%
% What each grid is there for:
%
%   GRIDMARK     three columns, two rows: the reported shape, and the one where
%                most tasks are not first in their row.
%   DISPLAYMARK  opens straight after a display.
%   COLUMNMARK   one column, so every task is first in its row. These never
%                nested, and were still loose paragraphs with the label after
%                the text.
%   RUNINMARK    no blank line before or after: the paragraph that follows has
%                to land in the same text unit as the list.
%   PARTSMARK    inside a list item, twice. The second opens the item, the one
%                place tasks takes a different path before its first row.
%   RICHMARK     a task holding a list, and a task of two paragraphs. The inner
%                list has to be tagged as a list again.
%   LABELMARK    a formula for a label (tasks typesets every label twice, and
%                the measuring copy must leave nothing in the tree), and a label
%                with no counter, which makes an unordered list.
%   GATEDMARK    a grid in a {solution} this build discards.
\documentclass{didactic}

\begin{document}

\section{Grids}

\begin{exercise}
	GRIDMARK. Factor each polynomial.
	\begin{tasks}(3)
		\task $x^2 - 1$
		\task $x^2 - 4$
		\task $x^2 - 9$
		\task $x^2 - 16$
		\task $x^2 - 25$
		\task $x^2 - 36$
	\end{tasks}
\end{exercise}

\begin{exercise}
	DISPLAYMARK. Evaluate the function
	\[
		f(x) = x^2 + 1
	\]
	\begin{tasks}(3)
		\task $f(1)$
		\task $f(2)$
		\task $f(5)$
	\end{tasks}
\end{exercise}

\begin{example}
	COLUMNMARK. Do the following for the function above.
	\begin{tasks}(1)
		\task Find the domain.
		\task Locate the intercepts.
	\end{tasks}

	\solution Read both off the graph.
\end{example}

\section{Around other structure}

RUNINMARK. A grid in running text, with no blank line on either side.
\begin{tasks}(2)
	\task alpha
	\task beta
\end{tasks}
The paragraph goes on after the grid.

\begin{exercise}
	\begin{parts}
		\item PARTSMARK. A grid inside a part.
		\begin{tasks}(2)
			\task gamma
			\task delta
		\end{tasks}
		\item
		\begin{tasks}(2)
			\task epsilon
			\task zeta
		\end{tasks}
	\end{parts}
\end{exercise}

\begin{exercise}
	RICHMARK. Tasks that hold more than a line.
	\begin{tasks}(2)
		\task A list:
		\begin{itemize}
			\item inner one
			\item inner two
		\end{itemize}
		\task First paragraph.

		Second paragraph.
	\end{tasks}
\end{exercise}

\begin{exercise}
	LABELMARK. Labels that are not the default.
	\begin{tasks}[label = $\alph*_0$, label-width = 2em](2)
		\task eta
		\task theta
	\end{tasks}
	\begin{tasks}[label = \textbullet](2)
		\task iota
		\task kappa
	\end{tasks}
\end{exercise}

\begin{exercise}
	GATEDMARK. Solve $x^2 - 3x + 2 = 0$.
\end{exercise}

\begin{solution}[1in]
	The two roots:
	\begin{tasks}(2)
		\task $x = 1$
		\task $x = 2$
	\end{tasks}
\end{solution}

\end{document}
